\section{Introduction}\label{introduction}

Mobile IoT devices and crowdsourced services have transformed service
composition from a static optimization problem into a dynamic,
sequential decision-making challenge. Traditional service composition
assumes fixed service locations---a user selects from a static pool of
available services. This assumption breaks when services themselves are
moving: pedestrians carrying devices in a shopping center, students
walking across campus, or workers moving through a factory floor. The composition that
worked seconds ago may fail moments later as services move out of
communication range.

\hypertarget{wifi-hotspot-sharing-scenario}{%
	\subsubsection{WiFi Hotspot Sharing
		Scenario}\label{wifi-hotspot-sharing-scenario}}

A representative example of moving crowdsourced services is WiFi hotspot
sharing through smartphones \cite{neiat2021, neiat2015}. This type of crowdsourced IoT service is
characterized by spatio-temporal aspects---the location/space and
time/period in which services are provisioned and consumed.

We identify two key types of crowdsourced services with regard to
spatial location: fixed and moving \cite{neiat2015}. A fixed crowdsourced service refers
to services that are permanent in space during the service provisioning
period---for example, a WiFi hotspot shared while sitting at a coffee
shop. In contrast, a moving crowdsourced service is not tied to any
specific location at any point in time---for example, sharing WiFi while
strolling through a shopping center or walking in the city.

Mobility presents key challenges for qualitative factors such as
availability. Fixed hotspot services remain available at a known
location when selected. However, for moving hotspot services, both
availability and location change during service provisioning.
Additionally, crowdsourced services may be deterministic (time period
and location are known in advance) or non-deterministic (unknown in
advance). This work focuses on deterministic moving services.

\hypertarget{research-challenges}{%
	\subsubsection{Research Challenges}\label{research-challenges}}

We identify three key research challenges for moving IoT service
composition.

The first challenge is connectivity---the core requirement for service
discovery. A moving service must stay connected with a consumer, meaning
it must remain within connectivity proximity. This requires determining
co-movement patterns between the user and service trajectories \cite{gudmundsson2006}. We
propose spatio-temporal filtering to find services that overlap with the
user trajectory in both space and time \cite{jeung2008}.

The second challenge is service continuity. A consumer and service
provider may not share their entire route---they may only overlap for
part of the journey. Therefore, an effective composition approach is
required to select an optimal sequence of available moving services that
ensure continuity \cite{neiat2015}. This is fundamentally different from static
composition where services remain available throughout \cite{deng2016}.

The third challenge is indexing and scalability. Existing co-movement
discovery methods rely on centralized index structures like R-trees. As
datasets scale up, performance degrades dramatically. We need an
approach that discovers and composes services without relying on
expensive indexing \cite{zhang2015}.

Moving IoT services exhibit spatio-temporal variability. Service
positions, availability, and quality attributes change continuously over
time \cite{deng2017}. A service might be within communication range at one instant and
outside it the next. Quality metrics like channel capacity depend on
distance, which changes as services move \cite{zhang2015}. This dynamic nature
fundamentally alters the composition problem from a one-time
optimization to a sequential decision process requiring real-time
adaptation \cite{gai2018}.

The core challenge is spatio-temporal validity. A service is valid only
when it falls within the consumer's discovery zone (communication range)
at the current time. Beyond validity, we want services that provide high
quality---the Signal Transmission Reward (STR) model derives capacity
from distance through exponential signal attenuation. The agent must
learn to select services that are both valid and high-quality, without
knowing a priori which services will be valid at each timestep.

Prior work addressed this with Double DQN \cite{neiat2021}, using the STR model
for distance-derived capacity and composition. The approach achieved
92.4\% success at low pedestrian mobility (2 km/h) and 71.8\% at higher
mobility (10 km/h), with frequent re-compositions needed \cite{neiat2021}. However,
DQN's value-based nature has three key limitations \cite{gai2018}. First,
overestimation bias leads to suboptimal action selection, particularly
harmful when valid services are few. Second, DQN handles discrete action
spaces poorly for service composition. Third, DQN is reactive---it
considers only current service positions, ignoring future states where
services may move out of range.

We propose Advantage Actor-Critic (A2C) to address these limitations.
A2C combines policy-based and value-based learning: the actor learns a
stochastic policy directly, while the critic estimates value. The
advantage function reduces variance while keeping updates unbiased,
enabling faster convergence and better sample efficiency. A2C naturally
handles discrete action spaces and can be extended with LSTM for
trajectory encoding to anticipate future states.

\textbf{Our Contributions}: This work makes four key contributions: (1) We adapt A2C for moving IoT service composition, building upon the STR-based selection from prior work \cite{neiat2021} while replacing Double DQN with A2C; (2) We implement and compare two network architectures; (3) We evaluate on real GPS datasets; (4) We analyze scalability with convergence within 350-500 episodes and 44% reduction in re-composition frequency.

We pose four research questions:

\begin{enumerate}
	\item How can A2C be adapted for moving IoT service composition while
	      building upon the STR-based selection from prior work \cite{neiat2021}?
	\item How do shared versus separate network architectures perform in this
	      domain?
	\item How does A2C compare to reactive baselines (greedy nearest-neighbor,
	      random selection)?
	\item How does A2C scale with increasing numbers of moving services (20 to
	      100+)?
\end{enumerate}

Section 2 covers background. Section 3 presents the A2C framework.
Section 4 describes experimental setup. Section 5 shows results. Section
6 provides implementation details. Section 7 discusses implications and
limitations. Section 8 concludes.
